% modulation.tex - clause 10 of v_cvm.tex: the signal set
\subsection{Signal set}\label{sec:modes}
The signal set has three modes, Table~\ref{tab:modes}.  A DCE shall
support the robust mode and the narrowband mode; it may support the
wideband mode (which needs the signal interface at 16000~Hz).

\begin{table}[htbp]
\caption{The modes of the signal set}
\label{tab:modes}
\small
\begin{tabularx}{\textwidth}{>{\raggedright\arraybackslash}p{2.3cm}>{\raggedright\arraybackslash}X>{\raggedright\arraybackslash}X>{\raggedright\arraybackslash}X}
\toprule
\rowcolor{ituhead}\thd{} & \thd{Robust (R)} & \thd{Narrowband (N)} & \thd{Wideband (W)} \\
\midrule
Signal & the FSK of \rec{V.21}: channel~1 (980/1180~Hz) from the calling DCE, channel~2 (1650/1850~Hz) from the answering DCE, 300~bit/s & a comb of tones at 200, 300 or 400~Hz spacing in the band 300 to 3400~Hz (16, 11 or 8 tones), DPSK & a comb of tones at 200, 300 or 400~Hz spacing in the band 300 to 6900~Hz (33, 23 or 17 tones), DPSK \\
Symbol & 3.33~ms (one bit) & 20~ms (a codec frame): a useful part of whole cycles of the spacing and a cyclic prefix (clause~\ref{sec:timing}); the 200~Hz comb also 22.5~ms & as N \\
Modulation & binary FSK & DQPSK, or 8-DPSK on tones with a low enough phase error & as N \\
Raw rate & 300~bit/s & 1600~bit/s (16 tones at 200~Hz, DQPSK, 20~ms), 1100 (11 at 300~Hz), 800 (8 at 400~Hz); half as much again with 8-DPSK & 3300~bit/s (33 tones at 200~Hz, DQPSK, 20~ms), 2300 (23 at 300~Hz), 1700 (17 at 400~Hz); half as much again with 8-DPSK \\
Use & INFO and MP messages, control sequences; data through codecs that carry nothing else & data & data \\
Passes (measured) & every AMR mode but 4.75, every AMR-WB mode but 6.6, every EVS mode but 9.6 wideband & AMR 12.2 (every comb), AMR 10.2 and two AMR 12.2 codecs in tandem (the 300 and 400~Hz combs), AMR-WB from 12.65, EVS 13.2 (N400) and 24.4 (every comb, 8-DPSK) & AMR-WB 23.85 (28 tones of W200), 15.85 (16 of W300), 12.65 (7 of W300); EVS 24.4 (W300), 48 and above (every comb, 8-DPSK); the super wideband combs through EVS at 48~\kbs{} and above \\
\bottomrule
\end{tabularx}
\end{table}

\itunote{The modes were chosen from measurements (Appendix~\ref{app:measurements}).
Through AMR at 7.4~\kbs{} and below, no differential phase modulation
measured kept a symbol error ratio below $10^{-1}$, on one carrier or many;
only 300~bit/s FSK did.  Through AMR at 12.2 and AMR-WB at 23.85~\kbs{},
a comb of 16 tones at 200~Hz spacing with DQPSK and a symbol of 20 to
22.5~ms kept its tone symbol error ratio at or below $10^{-2}$; combs of
32 tones at 100~Hz spacing did not (the codec's noise against each tone
grows faster than the number of tones), and 8-DPSK held only on the
better tones.  Fewer tones, more widely spaced, keep working where more
do not: through AMR at 10.2~\kbs{} and through two AMR codecs in tandem
the 11-tone comb at 300~Hz spacing kept $3 \cdot 10^{-3}$ while the
16-tone comb went to $3 \cdot 10^{-2}$, and through AMR-WB at 15.85 the
300~Hz comb loads 16 tones where the 200~Hz comb loads 10; so the signal
set has three combs for each band and the probing chooses.  A single
carrier at 600~baud (V.22's) passes AMR~12.2 as
well, at 1200~bit/s; the comb carries a third more at the same symbol
error ratio and loads itself to the codec.}

\subsection{The comb}\label{sec:comb}
The comb's tones are at the frequencies
\[ f_k = k \cdot \Delta f, \qquad k = k_{\min}, \ldots, k_{\max}, \]
as Table~\ref{tab:combs} gives them.
\begin{table}[htbp]
\caption{The combs}
\label{tab:combs}
\small
\begin{tabular}{llcccccc}
\toprule
\rowcolor{ituhead}\thd{Code} & \thd{Comb} & $\Delta f$ (Hz) & $k$ & \thd{Tones} & \thd{Band (Hz)} & $T_u$ (ms) & $T_g$ (ms) \\
\midrule
0 & N200 & 200 & 2--17 & 16 & 400--3400 & 15 (or 20) & 5 (or 2.5) \\
1 & N300 & 300 & 1--11 & 11 & 300--3300 & 10 & 10 \\
2 & N400 & 400 & 1--8 & 8 & 400--3200 & 15 & 5 \\
3 & W200 & 200 & 2--34 & 33 & 400--6800 & 15 & 5 \\
4 & W300 & 300 & 1--23 & 23 & 300--6900 & 10 & 10 \\
5 & W400 & 400 & 1--17 & 17 & 400--6800 & 15 & 5 \\
6 & S200 & 200 & 2--69 & 68 & 400--13800 & 15 & 5 \\
7 & S300 & 300 & 1--46 & 46 & 300--13800 & 10 & 10 \\
8 & S400 & 400 & 1--35 & 35 & 400--14000 & 15 & 5 \\
9--15 & -- & \multicolumn{6}{l}{reserved} \\
\bottomrule
\end{tabular}
\end{table}
The narrowband combs N200, N300 and N400 are for the signal interface
at 8000~Hz, the wideband combs W200, W300 and W400 for 16000~Hz, and the
super wideband combs S200, S300 and S400 for 32000~Hz (for the EVS
codec's super wideband modes, Appendix~\ref{app:evs}).  The probing (clause~\ref{sec:probing}) decides which
tones are loaded; the transmitter sends nothing on a tone that is not.
The transmitted signal in the useful part of symbol $n$ is
\[ s_n(t) = A \sum_{k \in \mathcal{L}} \cos\bigl(2\pi f_k t + \varphi_{n,k}\bigr), \qquad 0 \le t < T_u, \]
where $\mathcal{L}$ is the set of loaded tones, $\varphi_{n,k}$ the tone's
phase in this symbol and $A$ the amplitude that gives the level of
clause~\ref{sec:pointp}~b) over the loaded tones.  The useful period $T_u$
is a whole number of cycles of $\Delta f$, so that the tones are
orthogonal over it.

\subsection{Symbol timing and the cyclic prefix}\label{sec:timing}
A symbol consists of a cyclic prefix of $T_g$ followed by the useful part
of $T_u$: the prefix is the last $T_g$ of the useful part, transmitted
first, so that the signal is continuous across each symbol except for the
phase steps at the symbol's start.  The symbol periods are
\begin{itemize}
\item \textbf{20~ms}, the period of a codec frame, for every comb: $T_u$
  and $T_g$ as Table~\ref{tab:combs} gives them (three cycles of the
  spacing for the 200 and 300~Hz combs, six for the 400~Hz combs), so
  that one symbol's phase steps fall in each frame and the prefix spans
  the codec's transient after them (a subframe of 5~ms, or two for the
  300~Hz combs);
\item \textbf{22.5~ms}, for the N200 comb only: $T_u = 20$~ms (four
  cycles), $T_g = 2.5$~ms, for codecs that gain from the longer integration
  (the probing compares the two).
\end{itemize}
The receiver shall integrate each tone over the useful part only.  A
receiver shall acquire the symbol timing from the cyclic prefix (the
correlation of the signal with itself $T_u$ earlier peaks over the prefix)
and shall re-acquire it within two symbols after a step of the received
timing of up to $\pm 20$~ms, which a jitter buffer's adaptation causes
(Appendix~\ref{app:terminal}).  Differential modulation makes the symbol
after such a step the only one lost.

\subsection{The codec frame and the symbol}\label{sec:align}
The symbols are not aligned to the codec's frames, and no alignment is
signalled: the position of the symbols' phase steps within the frame was
found to make no difference to the tone symbol error ratio
(Appendix~\ref{app:measurements}, Table~\ref{tab:combnb}: the 20~ms
symbol at a frame boundary and 10~ms off it), which the codec's 5~ms
subframes and its long-term predictor explain: the step is coded in the
subframe it falls in, wherever that is.  The 20~ms symbol is chosen for
its rate (12 per cent more than the 22.5~ms symbol) and for its match
with the frame erasures: one erased frame is one symbol.

\subsection{Differential phase modulation}
Each loaded tone carries, per symbol, the phase difference
$\varphi_{n,k} - \varphi_{n-1,k}$, one of four values for DQPSK or eight for
8-DPSK, mapped from the bits as Table~\ref{tab:gray} shows.  The first
symbol of phase~4, and the first after a renegotiation, carries no data:
it is the phase reference.  The receiver decides each tone symbol from the
difference of the tone's phases in two consecutive symbols; no carrier
recovery and no equalization are required, and a frequency offset between
the DCEs' clocks of $\Delta f_c$ appears as a constant phase rotation of
$2\pi \Delta f_c T$ per symbol, which the receiver should estimate from
the decisions and remove.

\subsection{Loading and raw rate}
The raw rate is
\[ R = \frac{1}{T_u + T_g} \sum_{k \in \mathcal{L}} b_k, \]
with $b_k = 2$ for a DQPSK tone and 3 for an 8-DPSK tone: 1600~bit/s for the
16 DQPSK tones of N200 at the 20~ms symbol, 1100 for N300, 800 for N400,
3300 for W200, and half as much again with 8-DPSK on all tones.  The loading of clause~\ref{sec:loading} chooses $\mathcal{L}$
and the $b_k$; a loading with fewer than 4 tones shall not be used (the
DCE renegotiates or retrains).

\subsection{Control sequences during data}\label{sec:control}
A DCE that wishes to renegotiate (clause~\ref{sec:reneg}) shall send, in
place of data, the control sequence: eight consecutive symbols with the
phase difference $180^\circ$ on every loaded tone, then its MP in the
robust mode.  A receiver that decodes four such symbols in a row shall
stop delivering data to the inner decoder, flush, answer with the same
sequence if it has not sent one, and proceed to the MP exchange.  The
sequence cannot occur in scrambled data for more than two symbols with a
probability above $10^{-9}$ on 8 tones.

\subsection{Scrambling}
The bits given to the mapping of clause~\ref{sec:mapping} shall be
scrambled with the generating polynomials of \rec{V.34}: $1 + x^{-18} + x^{-23}$
in the calling DCE's transmitter and $1 + x^{-5} + x^{-23}$ in the
answering DCE's (GPC and GPA), self-synchronizing as in that
Recommendation.  The scrambler is reset at the start of phase~4 and after
each renegotiation.

\subsection{Quick start}\label{sec:quick}
A DCE may keep, for the last number dialled or answered, the loading and
the symbol period of the last call, and offer them in INFO0 as the mode
to start with.  If its peer accepts in its INFO0, phase~2 is reduced to L1 of
0.5~s for a check of the loading, and the DCEs proceed to phase~3; this
is in the manner of \rec{V.92}'s quick connect.  A loading so offered
shall still be checked by the receiver against the thresholds of
clause~\ref{sec:loading}.

\subsection{Loading with amplitude: for further study}\label{sec:qam}
The MP message carries a loading value reserved for QAM constellations of
16 and 64~points per tone, with the 16-state four-dimensional trellis code
of \rec{V.34} across tone pairs as \rec{G.992.1} applies it.  Such a
loading is for the transform coding modes of the EVS codec, which carry
waveforms: through EVS at 48~\kbs{} and above a tone's amplitude was
measured to change by 0.3 to 1~dB rms from one symbol to the next, which
would hold two levels 6~dB apart (Appendix~\ref{app:evs}); through the
ACELP codecs it changes by 1.5 to 3~dB, and amplitude carries no bits.
The loading's parameters are for further study.  A DCE shall not
request it unless both DCEs have indicated the capability in CM and JM,
and shall fall back to DPSK when the peer does not confirm it in MP.
